\documentclass[11pt,a4paper]{article}

\usepackage[T1]{fontenc}
\usepackage[utf8]{inputenc}
\usepackage[main=italian,provide=*]{babel}
\usepackage{amsmath,amssymb,mathtools}
\usepackage{booktabs}
\usepackage{geometry}
\geometry{margin=2.6cm}
\usepackage{xcolor}
\usepackage{fancyhdr}
\usepackage{enumitem}
\usepackage{seqsplit}
\usepackage{hyperref}
\hypersetup{colorlinks=true,linkcolor=blue,citecolor=blue,urlcolor=blue}

\newcommand{\ket}[1]{\lvert #1\rangle}
\newcommand{\vs}[1]{\mathbf{s}_{#1}}

\pagestyle{fancy}
\fancyhf{}
\fancyhead[L]{\small Trimero a catena aperta — analisi del termine DM}
\fancyhead[R]{\small\thepage}
\renewcommand{\headrulewidth}{0.3pt}

\title{\bfseries Il termine DM sulla catena aperta:\\
perché apre il gap pur avendo la stessa struttura dell'Opzione A dell'anello}
\author{Note di lavoro — Tesi triennale, Samuele Schiavi\\ \small Relatore: Prof.\ Alessandro Chiesa}
\date{}

\begin{document}
\maketitle
\thispagestyle{fancy}

\begin{abstract}
\noindent
Questo documento risponde al punto 1 della Fase 5 per la catena aperta: la struttura
del termine DM è già nota ($D_{12}=-D_{23}\equiv D$, unica compatibile con la
riflessione $P_{13}$ — si veda \texttt{teoria\_trimero\_catena\_aperta.pdf}), qui va
solo \emph{quantificata}. Il risultato non è un mirror banale dell'anello: il DM della
catena conserva $S_{13}^2$ esattamente, la stessa struttura dell'Opzione A
dell'anello — che non apriva \emph{mai} il gap. Sulla catena il gap si apre invece
linearmente, $g_\text{min}(D)=2\sqrt6\,D$. La spiegazione è la stessa regola di
simmetria usata per l'anello, applicata al caso complementare.
\end{abstract}

\section{Il punto di partenza}

Nella teoria esatta della catena (decomposizione di Kambe rispetto alla coppia
\emph{non legata} $(1,3)$), l'Hamiltoniana
\[
H_0 = J(\boldsymbol\sigma_1\!\cdot\!\boldsymbol\sigma_2+\boldsymbol\sigma_2\!\cdot\!\boldsymbol\sigma_3)
+ b\sum_{i=1}^3 Z_i
\]
si risolve in forma chiusa in tre multipletti: $A'$ ($S_{13}{=}1,S{=}\tfrac32$, $E_A=2J$),
$B'$ ($S_{13}{=}1,S{=}\tfrac12$, $E_B=-4J$), $C'$ ($S_{13}{=}0,S{=}\tfrac12$, $E_C=0$).
Per $J>0$ il fondamentale a basso campo è \emph{sempre} $B'$ (mai serve confrontare
con $C'$: $E_B<E_C$ indipendentemente da $J$), con incrocio vero a $b_c=3J$ verso
$A'$ — dimostrato esatto (gap nullo, non anticrossing) perché $H_0$ è diagonale nei
numeri quantici $(S_{13},S,M)$.

\textbf{La domanda di questo documento.} Il termine DM compatibile con $P_{13}$,
\[
H_{DM} = D\,(X_1Z_2-Z_1X_2) - D\,(X_2Z_3-Z_2X_3),
\]
apre questo incrocio? La struttura è nota da subito (\texttt{trimer\_chain\_exact.py},
self-test DM): $[P_{13},H_{DM}]=0$ e, conseguenza più profonda,
$[S_{13}^2,H_{DM}]=0$ esattamente. Quest'ultimo fatto è precisamente la struttura
dell'\textbf{Opzione A} dell'anello — quella che \emph{non apre mai} il gap
(\texttt{analisi\_dm\_trimero\_anello.tex}). Un mirror ingenuo concluderebbe che
anche qui il gap resta chiuso. È falso, e il motivo è istruttivo.

\section{Elemento di matrice all'incrocio}

L'incrocio a $b_c=3J$ è fra $\ket{A',M{=}-\tfrac32}=\ket{111}$ ed
$\ket{B',M{=}-\tfrac12}=\tfrac{1}{\sqrt6}(\ket{011}+\ket{110})-\sqrt{\tfrac23}\ket{101}$.
Applicando $H_{DM}$ a $\ket{111}$:
\[
H_{DM}\ket{111} = D\big({-}\ket{011}+2\ket{101}-\ket{110}\big),
\]
da cui, proiettando su $\ket{B',-\tfrac12}$:
\[
\big\langle B',{-}\tfrac12\big|H_{DM}\big|A',{-}\tfrac32\big\rangle = -\sqrt6\,D.
\]
Verificato numericamente contro \texttt{trimer\_chain\_exact.py::kambe\_states()}:
errore $4.4\times10^{-16}$. In teoria delle perturbazioni degeneri (i due stati sono
esattamente degeneri a $D=0$, $b=b_c$), questo dà
\[
g_\text{min}(D) \simeq 2\,|V| = 2\sqrt6\,D \approx 4.899\,D \qquad (D\ \text{piccolo}).
\]

\section{Verifica numerica: \texorpdfstring{$g_\text{min}(D)$}{g\_min(D)}}
\label{sec:gmin-tabella}

Calcolato $g_\text{min}(D)=\min_{b\in[b_c-w,b_c+w]}[E_1(b;D)-E_0(b;D)]$, con la
stessa metodologia a due cautele già stabilita per l'anello
(\texttt{analisi\_dm\_trimero\_anello.tex}, sez.\ ``Il controllo corretto'' e sez.\
``Una seconda trappola''): mai il gap a $b_c$ fisso, mai il minimo su tutto l'asse
$b$ (a $b=0$ la catena ha la stessa degenerazione di Kramers dell'anello — verificato
esplicitamente, $[T,H]=0$ a $b=0$ per ogni $D$, $T^2=-\mathbb I$, tre spin-$1/2$).

\begin{center}
\renewcommand{\arraystretch}{1.3}
\begin{tabular}{cccc}
\toprule
$D$ & $g_\text{min}$ & $b$ al minimo & $g_\text{min}/D$ \\
\midrule
$0.009$ & $0.044091$ & $3.0000068$ & $4.8990$ \\
$0.020$ & $0.097979$ & $3.0000333$ & $4.8990$ \\
$0.050$ & $0.244945$ & $3.0002082$ & $4.8989$ \\
$0.073$ & $0.357612$ & $3.0004435$ & $4.8988$ \\
$0.100$ & $0.489864$ & $3.0008314$ & $4.8986$ \\
$0.148$ & $0.724939$ & $3.0018159$ & $4.8982$ \\
$0.200$ & $0.979524$ & $3.0033019$ & $4.8976$ \\
$0.300$ & $1.468777$ & $3.0073414$ & $4.8959$ \\
$0.500$ & $2.445262$ & $3.0196198$ & $4.8905$ \\
\bottomrule
\end{tabular}
\end{center}

Gap lineare in $D$ con pendenza $2\sqrt6=4.898979$, confermata dalla predizione
perturbativa. \textbf{Differenza qualitativa dall'anello}: qui il punto critico si
sposta pochissimo ($b_\text{min}$ da $3.0000068$ a $3.0196$ su tutto il range),
contro $2.40\to2.43$ (Opzione B) o $2.40\to2.29$ (Opzione A) dell'anello — il DM
accoppia direttamente i due stati quasi degeneri senza spostare i rami al primo
ordine.

\section{Perché si apre: la stessa regola, il caso complementare}

La regola di non-incrocio di von Neumann--Wigner (già usata per l'anello): un
termine che preserva un numero quantico non può aprire un incrocio fra stati che
quel numero quantico \emph{distingue}. Il punto decisivo è quali stati si
incrociano:

\begin{center}
\renewcommand{\arraystretch}{1.4}
\begin{tabular}{@{}p{3.4cm}p{5.6cm}p{5.6cm}@{}}
\toprule
& \textbf{Anello, Opzione A} & \textbf{Catena} \\
\midrule
incrocio & $C$ ($S_{12}{=}0$) $\to$ $A$ ($S_{12}{=}1$) & $B'$ ($S_{13}{=}1$) $\to$
$A'$ ($S_{13}{=}1$) \\
settori del Casimir conservato & \textbf{diversi} & \textbf{uguali} \\
elemento di matrice & vincolato a zero per costruzione & \emph{non} vincolato \\
esito & incrocio vero per ogni $D$ & anticrossing, $g_\text{min}=2\sqrt6\,D$ \\
\bottomrule
\end{tabular}
\end{center}

Il DM della catena conserva $S_{13}^2$ esattamente — verificato,
$\|[H+H_{DM},S_{13}^2]\|=0$ su 20 punti casuali — ma l'incrocio a $b_c$ è
\emph{intra-settore}: entrambi gli stati hanno $S_{13}=1$. Il Casimir conservato non
li distingue, quindi non li protegge. Controllo diretto sul fondamentale a $b=b_c$:
$\langle S_{13}^2\rangle=2.000000$ per entrambi i due livelli più bassi (qualunque
$D$), mentre $\langle S^2\rangle$ e $\langle M_z\rangle$ sono mescolati.

\begin{center}
\renewcommand{\arraystretch}{1.3}
\begin{tabular}{ccccc}
\toprule
$D$ & $|\langle A'|gs\rangle|^2$ & $|\langle B'|gs\rangle|^2$ & $\langle M_z\rangle$ &
gap \\
\midrule
$0.005$ & $0.499915$ & $0.500085$ & $-0.9999$ & $0.024495$ \\
$0.150$ & $0.497451$ & $0.502401$ & $-0.9973$ & $0.734742$ \\
$0.500$ & $0.491583$ & $0.506972$ & $-0.9901$ & $2.445579$ \\
\bottomrule
\end{tabular}
\end{center}

A $b=b_c$ la miscela è $50/50$ per \emph{qualunque} $D>0$ (i due stati sono
degeneri a $D=0$): $D$ fissa il gap, non il mescolamento — la rottura di $M$ è
massimale già a $D$ infinitesimo. Fuori da $b_c$ il fondamentale torna rapidamente
quasi puro ($|\langle A'|gs\rangle|^2=0.031$ a $b=2.0$, $D=0.15$): $b=b_c$ è il
punto di massima severità per un ansatz $M$-conservante, ed è per questo il punto
scelto per la Fase 5 (VQE).

\section{Punto di lavoro adottato}

\[
\boxed{J=1,\qquad b=b_c=3J=3.0,\qquad D=0.15}
\]

Stesso $D$ già usato per l'anello (confronto diretto a parità di ampiezza DM, non
di gap). Gap risultante $0.7347$: ben risolto, fondamentale non degenere, primo
eccitato isolato ($E_2-E_1=3.63$). Ancora perturbativo: peso del fondamentale sul
sottospazio degenere di $D=0$ pari a $0.999852$.

\section{Nota metodologica: assunzione presa in autonomia}

Come per il Trotter dell'anello (Opzione B assunta senza conferma esplicita), la
quantificazione del DM per la catena non era stata richiesta esplicitamente dal
relatore (si veda \texttt{domande\_relatore.md}) — svolta per coerenza con il lavoro
già fatto sull'anello e per completare il parallelismo fra le due topologie. Da
segnalare nel prossimo report.

\section{Conclusioni}

\begin{enumerate}[leftmargin=1.4em]
\item Il DM della catena apre l'incrocio $B'\to A'$ a $b_c=3J$, con
$g_\text{min}(D)=2\sqrt6\,D$ verificato su griglia.
\item L'analogia strutturale con l'Opzione A dell'anello (stesso Casimir
conservato) è \emph{fuorviante se letta come analogia di esito}: l'anello con
Opzione A non apre mai il gap, la catena sì, perché l'incrocio è fra stati con lo
stesso numero quantico $S_{13}$ (o $S_{12}$), non fra settori diversi.
\item Punto di lavoro $(J,b,D)=(1,3.0,0.15)$ confermato, stessa severità di
mescolamento dell'anello a $b_c$ ma raggiunta con un solo grado di libertà DM
anziché due opzioni da confrontare.
\end{enumerate}

\end{document}
